% Bosquejo del método (metodo.tex), según el material «Bosquejo del método»
% de Taller de Investigación I (Dra. Morales Rosales). Sigue sus seis
% apartados: universo y muestra, tipo de estudio, instrumento, plan de
% recolección, plan de procesamiento y análisis y plan de presentación
% gráfica. Retoma las decisiones de secciones/05-metodologia.tex y de los
% Apéndices A a D; si cambia alguna, actualizar ambos archivos.
% Las fechas de la Tabla tab:calendario-metodo son una propuesta que el
% equipo debe confirmar con el calendario escolar y las autorizaciones.
\section{Bosquejo del Método}
El bosquejo del método describe los procedimientos con que se responderá la
pregunta de investigación: a quiénes se estudiará, qué tipo de investigación
se realizará, qué información se necesita y con qué instrumentos se
obtendrá, cuándo y cómo se recopilará, cómo se organizarán e interpretarán
los datos y cómo se presentarán los resultados.

\subsection{Universo y Obtención de la Muestra}
El universo está formado por los estudiantes inscritos en el \tescha{} que
ingresan a pie al plantel. La investigación no los estudia a todos por igual,
porque cada instrumento necesita un subconjunto distinto de esa población y
de las personas que intervienen en el acceso. Por eso se definen tres
poblaciones de estudio, cada una con su propia forma de selección.

La primera población son los ingresos al acceso peatonal piloto en el turno
vespertino, entre 30 minutos antes y 10 minutos después del inicio de la
primera clase. En ese periodo se cuentan todos los ingresos y se registra la
longitud máxima de la fila de cada intervalo de cinco minutos, de modo que
para estos dos datos no hay muestreo, sino un conteo completo. El tiempo de
validación, en cambio, no puede cronometrarse a cada estudiante, por lo que en
cada intervalo se cronometra a cinco estudiantes elegidos al azar, como indica
el Apéndice~C. Se trata de un muestreo probabilístico aleatorio simple dentro
de cada intervalo, en el que cada estudiante que llega tiene la misma
probabilidad de ser elegido \parencite{hernandezsampieri2014metodologia}. Con ocho intervalos por día durante cinco días hábiles, cada
serie de observación reúne 40 intervalos y 200 tiempos de validación.

La segunda población son los estudiantes de todas las carreras y turnos, a
quienes se dirige la encuesta del Apéndice~A. Como no se cuenta con una lista
de la matrícula a partir de la cual sortear a los participantes, el muestreo
es no probabilístico por cuotas: se fija un número de cuestionarios por turno
en proporción a la matrícula de cada uno, que proporcionará Control Escolar, y
los cuestionarios se aplican a quienes acepten responder hasta completar cada
cuota. La meta es reunir al menos 196 cuestionarios, el tamaño que la fórmula
del Apéndice~A da para una matrícula grande con un margen de error de 0.07.
Esa cifra se usa solo como referencia de tamaño: al no ser una muestra
probabilística, los resultados describen a quienes respondieron y no se les
atribuye un margen de error.

La tercera población es el personal que interviene en el acceso: los
vigilantes del acceso piloto y el personal del Departamento de Control
Escolar. Son pocas personas y se eligen por la información que poseen, por lo
que el muestreo es no probabilístico intencional, en el que la elección
depende del criterio del investigador y no de la probabilidad
\parencite{hernandezsampieri2014metodologia}. Se entrevistará al menos a
un integrante de cada área (Apéndice~B).

Además de las personas, el estudio somete a prueba el propio prototipo. Las
pruebas de banco no trabajan con una muestra de estudiantes, sino con un
número fijo de lecturas en cada condición de iluminación, brillo de pantalla
y distancia, y con una lista cerrada de casos de validación (Apéndice~E).

\subsection{Determinación del Tipo de Estudio}
La clasificación del estudio se hace desde varios criterios que no se
excluyen entre sí. La Tabla~\ref{tab:tipo-estudio} resume la decisión
tomada en cada uno y la razón que la sustenta.

\begin{table}[tbp]
	\caption{Tipo de Estudio Según Cada Criterio de Clasificación}
	\label{tab:tipo-estudio}
	\small\linespread{1}\selectfont
	\begin{tabular}{@{}P{2.3cm}P{3.6cm}P{9.6cm}@{}}
		\toprule
		Criterio     & Clasificación                                   & Razón                                                                                                                                                                                                         \\
		\midrule
		Finalidad    & Aplicada                                        & Usa conocimiento existente sobre códigos QR y teoría de colas para resolver un problema concreto del acceso, no para generar teoría nueva \parencite{oecd2015frascati}.                                      \\[0.5em]
		Enfoque      & Mixto, con predominancia cuantitativa           & Los indicadores principales se miden y se comparan numéricamente; la entrevista aporta información cualitativa sobre el procedimiento y las condiciones institucionales \parencite{hernandezsampieri2014metodologia}. \\[0.5em]
		Alcance      & Diagnóstico, descriptivo y explicativo          & Primero se diagnostica y describe el acceso actual para obtener la línea base; después se busca establecer si el cambio de mecanismo de validación se asocia con un cambio en los indicadores.             \\[0.5em]
		Fuente       & Documental y de campo                           & La fundamentación y los requisitos legales provienen de fuentes publicadas; los datos que responden la pregunta no existen en ningún registro y se levantan en el acceso.                                     \\[0.5em]
		Control      & Experimental, de tipo cuasiexperimental         & Se manipula el mecanismo de validación, pero sin asignación aleatoria ni grupo de control; el diseño es de preprueba y posprueba con un solo grupo \parencite{campbellstanley1963}.                      \\[0.5em]
		Temporalidad & Longitudinal                                    & Se mide en dos series separadas por la intervención, antes del piloto y durante él, sobre el mismo acceso y con la misma hoja de registro.                                                                   \\
		\bottomrule
	\end{tabular}
	\tablenote{Elaboración propia con los criterios de clasificación vistos en
		clase.}
\end{table}

En consecuencia, se trata de una investigación aplicada, de campo y
cuasiexperimental. Con la notación de \textcite{campbellstanley1963}, el
diseño se representa como $G\;O_1\;X\;O_2$, donde $G$ es el acceso piloto,
$O_1$ la primera serie de observación, $X$ la instalación del prototipo y
$O_2$ la segunda serie. Sus resultados se interpretarán como evidencia
orientativa para el acceso estudiado y no como una generalización, porque el
diseño no descarta por completo otras causas de cambio, como la historia y la
maduración \parencite{cranmer2017pretestposttest}. A ellas se suma la
reacción de las personas al saberse observadas o ante una novedad, conocida
como efecto Hawthorne \parencite{mccambridge2014hawthorne}. Para atenuarla,
los observadores se colocan fuera del paso, se dejan dos días de adaptación
al prototipo antes de la segunda serie y se registra el número de puntos de
revisión activos en cada jornada.

Las variables del estudio se definen de forma operacional en la
Tabla~\ref{tab:variables-metodo}, es decir, por la manera en que se medirán.

\begin{table}[tbp]
	\caption{Definición Operacional de las Variables}
	\label{tab:variables-metodo}
	\small\linespread{1}\selectfont
	\begin{tabular}{@{}P{3.2cm}P{6.2cm}P{2.6cm}P{3.1cm}@{}}
		\toprule
		Variable & Definición operacional & Unidad & Instrumento \\
		\midrule
		Mecanismo de validación (independiente) & Revisión visual de la credencial en la línea base; lectura del código QR dinámico durante el piloto. & Dos condiciones & Calendario del estudio \\[0.5em]
		Tiempo de validación (dependiente) & Tiempo desde que el estudiante llega al punto de revisión hasta que lo cruza; se promedia por intervalo. & Segundos & Apéndice~C \\[0.5em]
		Longitud máxima de la fila (dependiente) & Mayor número de personas formadas detrás del punto de revisión durante cada intervalo de cinco minutos. & Personas & Apéndice~C \\[0.5em]
		Incidencias de identidad (complementaria) & Ingresos sin credencial, sin revisión y rechazados, contados por separado en cada intervalo. & Casos & Apéndice~C \\[0.5em]
		Tasa de llegadas y utilización (de control) & Ingresos por minuto en cada intervalo y utilización $\rho = \lambda/(c\,\mu)$ con los puntos de revisión activos. & Estudiantes por minuto; sin unidad & Apéndice~C \\
		\bottomrule
	\end{tabular}
	\tablenote{Elaboración propia a partir de la hipótesis y del Apéndice~C.}
\end{table}

\subsection{Selección, Diseño y Prueba del Instrumento}
Cada técnica se eligió por el tipo de información que se necesita, y cada una
tiene su instrumento y su medio de aplicación
(Tabla~\ref{tab:tecnicas-metodo}).

\begin{table}[tbp]
	\caption{Técnicas, Instrumentos y Medios de Aplicación}
	\label{tab:tecnicas-metodo}
	\small\linespread{1}\selectfont
	\begin{tabular}{@{}P{3cm}P{3.4cm}P{3.4cm}P{5.3cm}@{}}
		\toprule
		Técnica                                         & Instrumento                                     & Medio de aplicación                                         & Información que aporta                                                                             \\
		\midrule
		Observación estructurada no participante        & Hoja de registro (Apéndice~C)                   & Hoja impresa y cronómetro del celular                       & Ingresos, longitud máxima de la fila, tiempo de validación e incidencias por intervalo.          \\[0.5em]
		Encuesta                                        & Cuestionario (Apéndice~A)                       & Google Forms y, si hace falta, copia impresa                & Percepción del acceso, disponibilidad de celular, conexión y batería, y disposición a usar el QR. \\[0.5em]
		Entrevista semiestructurada                     & Guía de entrevista (Apéndice~B)                 & Presencial, con notas escritas                              & Procedimiento real de revisión, matrícula por turno y condiciones para tratar datos personales.  \\[0.5em]
		Medición técnica                                & Protocolo y hoja de pruebas de banco (Apéndice~E) & Prototipo, cronómetro y registro del servidor               & Tiempo de una lectura completa y proporción de lecturas exitosas antes del piloto.               \\[0.5em]
		Revisión documental                             & Ficha de requisitos legales (Apéndice~F)        & Consulta de las leyes de protección de datos                & Requisitos del aviso de privacidad, del resguardo y del periodo de conservación.                 \\
		\bottomrule
	\end{tabular}
	\tablenote{Elaboración propia. La revisión documental se basa en la ley
		estatal \parencite{edomex2017datos} y en la ley general
		\parencite{lgpdppso2025} de protección de datos personales.}
\end{table}

Antes de aplicarlos, los instrumentos se probarán para corregir lo que no
funcione. El cuestionario se aplicará a un grupo de diez estudiantes que no
formarán parte de la muestra. Un grupo pequeño basta, porque la prueba
piloto de un instrumento busca comprobar su claridad y su factibilidad, no
estimar parámetros de la población \parencite{hertzog2008pilot}. Con esa
prueba se verificará que las preguntas se entiendan,
que las opciones de respuesta no se traslapen y que el tiempo de respuesta
sea de unos cinco minutos. La hoja de observación se ensayará durante un día
en el acceso piloto, con los dos observadores, para confirmar que es posible
contar los ingresos, anotar la fila y cronometrar a cinco estudiantes en cada
intervalo de cinco minutos; los datos de ese día no se suman a la línea base.
La guía de entrevista se revisará entre los integrantes del equipo para
ajustar el orden y la redacción de las preguntas. Los cambios que resulten de
estas pruebas se harán en los apéndices antes de la primera aplicación.

\subsection{Plan de Recolección de Información}
La información se obtendrá en el \tescha{}, en Chalco, Estado de México. La
observación se hará en el acceso peatonal piloto, en el turno vespertino y en
la franja que va de 30 minutos antes a 10 minutos después del inicio de la
primera clase. La encuesta se aplicará a estudiantes de ambos turnos y las
entrevistas se harán en el lugar de trabajo de cada informante. La
Tabla~\ref{tab:calendario-metodo} presenta el periodo propuesto para cada
actividad.

\begin{table}[tbp]
	\caption{Calendario Propuesto de Recolección de Información}
	\label{tab:calendario-metodo}
	\small\linespread{1}\selectfont
	\begin{tabular}{@{}P{6.2cm}P{4.4cm}P{4.9cm}@{}}
		\toprule
		Actividad                                                             & Periodo                                   & Instrumento o producto                          \\
		\midrule
		Solicitud de autorización a la Dirección y a Control Escolar          & 18 al 20 de noviembre de 2026             & Oficio de solicitud                             \\[0.5em]
		Revisión de la legislación y ficha de requisitos legales             & 18 al 27 de noviembre de 2026             & Ficha de requisitos (Apéndice~F)                \\[0.5em]
		Prueba del cuestionario y ensayo de la hoja de observación            & 23 y 24 de noviembre de 2026              & Apéndices A y C corregidos                      \\[0.5em]
		Entrevistas a vigilancia y a Control Escolar                          & 23 al 27 de noviembre de 2026             & Guía de entrevista (Apéndice~B)                 \\[0.5em]
		Aplicación de la encuesta                                             & 25 de noviembre al 4 de diciembre de 2026 & Cuestionario (Apéndice~A)                       \\[0.5em]
		Pruebas de banco del prototipo                                        & 30 de noviembre al 4 de diciembre de 2026 & Hoja de pruebas de banco (Apéndice~E)          \\[0.5em]
		Primera serie de observación (línea base)                             & 30 de noviembre al 4 de diciembre de 2026 & Hoja de observación (Apéndice~C)                \\[0.5em]
		Instalación, alta de estudiantes y adaptación al prototipo            & 7 y 8 de diciembre de 2026                & Aviso de privacidad y consentimiento (Apéndice~G) \\[0.5em]
		Segunda serie de observación (piloto)                                 & 9 al 15 de diciembre de 2026              & Hoja de observación (Apéndice~C)                \\[0.5em]
		Procesamiento y análisis de los datos                                 & 16 de diciembre de 2026 al 15 de enero de 2027 & Base de datos y resultados                 \\[0.5em]
		Presentación de resultados                                            & 18 al 29 de enero de 2027                 & Informe final                                   \\
		\bottomrule
	\end{tabular}
	\tablenote{Elaboración propia a partir del calendario del plan de gestión
		del proyecto, que fija las pruebas y el piloto hasta el 18 de diciembre.
		Las fechas están sujetas al calendario escolar, a la autorización del
		plantel y a que el prototipo pase las pruebas de banco. Las dos series de
		observación duran cinco días hábiles cada una; la del piloto abarca del
		9 al 11 y del 14 al 15 de diciembre. El receso del 21 de diciembre al 1
		de enero no se cuenta.}
\end{table}

El procedimiento de recolección seguirá estos pasos:
\begin{enumerate}
	\item Solicitar por escrito la autorización de la Dirección del plantel y
	      de Control Escolar para observar el acceso, aplicar la encuesta y
	      realizar el piloto.
	\item Informar a los participantes el propósito del estudio: a los
	      estudiantes, en el aviso de privacidad del cuestionario; al personal,
	      al inicio de la entrevista.
	\item Obtener el consentimiento: la respuesta voluntaria al cuestionario,
	      el consentimiento verbal de los entrevistados y la autorización escrita
	      de los estudiantes que se den de alta en el piloto, con el formato del
	      Apéndice~G.
	\item Aplicar la encuesta y las entrevistas, y levantar la primera serie de
	      observación durante cinco días hábiles consecutivos.
	\item Instalar el prototipo, dar de alta a los estudiantes participantes y
	      dejar dos días de adaptación antes de medir.
	\item Levantar la segunda serie de observación durante cinco días hábiles,
	      con el mismo acceso, la misma franja, los mismos observadores y la
	      misma hoja, y anotar en cada hoja el número de puntos de revisión
	      activos.
	\item Verificar al final de cada jornada que las hojas estén completas,
	      capturarlas el mismo día y resguardar los originales.
\end{enumerate}

Los seis integrantes del equipo comparten la recolección. Dos de ellos
actúan como observadores en las dos series, para que las mediciones sean
comparables; el primero cuenta los ingresos y la fila, y el segundo
cronometra. Los demás aplican la encuesta, realizan las entrevistas, hacen
las pruebas de banco y capturan los datos. Los recursos necesarios son las
hojas de observación impresas, cronómetros (los de los teléfonos de los
observadores), el formulario en línea y copias impresas del cuestionario,
la guía de entrevista y el prototipo con su servidor y su panel web.

En cuanto a las consideraciones éticas, la encuesta es anónima y voluntaria y
no pide nombre ni matrícula; la entrevista no registra nombres; la
observación anota conteos y tiempos, nunca datos que identifiquen a un
estudiante. Los datos del piloto se tratarán conforme a la legislación de
protección de datos personales \parencite{edomex2017datos, lgpdppso2025},
según los requisitos reunidos en el Apéndice~F: los estudiantes recibirán el
aviso de privacidad del Apéndice~G antes de darse de alta, solo se
registrarán los datos necesarios para validar el acceso, y los registros se
conservarán durante el periodo que se acuerde con Control Escolar. Quien no
quiera participar en el piloto seguirá entrando con la revisión manual.

\subsection{Plan de Procesamiento y Análisis}
Las hojas de observación se capturarán en una hoja de cálculo con un renglón
por intervalo, y las respuestas del cuestionario se descargarán del
formulario en línea. Antes del análisis se revisará que no falten datos, se
eliminarán los cuestionarios duplicados o incompletos, se corregirán las
inconsistencias con la hoja original y se codificarán las respuestas: las
opciones de la escala de Likert con valores del 1 al 5 y las demás preguntas
cerradas como categorías. Las respuestas de la escala son datos ordinales: la
distancia entre «de acuerdo» y «totalmente de acuerdo» no es necesariamente
igual a la que hay entre otras dos opciones, por lo que se resumirán con la
mediana y las frecuencias y no con la media \parencite{jamieson2004likert}.
Los cálculos se harán con scripts de Python, que se
conservarán junto con los datos para que el análisis pueda repetirse.

El análisis cuantitativo tiene dos partes. La primera es descriptiva: con los
datos de cada intervalo se calculan la tasa de llegadas, el tiempo medio de
validación, la tasa de servicio y la utilización del acceso, como indica el
Apéndice~C, y se identifican los intervalos con utilización igual o mayor
que 1. Para cada serie se informan la media, la mediana y la desviación
estándar del tiempo de validación y de la longitud máxima de la fila. De la
encuesta se obtienen frecuencias y porcentajes de cada respuesta, y la
mediana de cada afirmación de la escala. La viabilidad tecnológica se
expresa como la proporción de estudiantes que dicen llevar siempre o casi
siempre un celular propio, con conexión y con batería al llegar; la
disposición a usar el sistema, como la proporción que marca 4 o 5 en la
afirmación 13.

La segunda parte contrasta la hipótesis. Los estudiantes cronometrados antes
y durante el piloto no son los mismos, de modo que la unidad de análisis es
el intervalo: cada uno de los 40 intervalos de la línea base se empareja con
el intervalo que ocupa la misma posición, en el mismo día de la serie, en el
piloto. Con esos 40 pares se compara el tiempo medio de validación y la
longitud máxima de la fila. Primero se revisará la normalidad de las
diferencias con la prueba de Shapiro-Wilk \parencite{shapiro1965normality}.
Si las diferencias tienen una distribución aproximadamente normal, se usará
la prueba \emph{t} para muestras relacionadas; si no, la prueba de rangos con
signo de Wilcoxon \parencite{wilcoxon1945ranking}, que no supone normalidad.
En ambos casos se usará un nivel de significancia de 0.05 y la hipótesis
nula de que el prototipo no reduce el tiempo de validación ni la fila. Se
informará, además, el tamaño del efecto y la magnitud de la diferencia en
segundos y en personas, porque el valor \emph{p} indica si hay una
diferencia, pero no qué tan grande es \parencite{sullivan2012effectsize}.

Las incidencias de identidad, es decir, los ingresos sin credencial, sin
revisión y rechazados, se resumirán por separado con frecuencias por serie,
sin prueba estadística, como información complementaria sobre la seguridad.
Las pruebas de banco se resumirán con el tiempo medio de una lectura completa
y la proporción de lecturas exitosas en cada condición de iluminación,
brillo y distancia, y se compararán con los criterios de aceptación del
Apéndice~E antes de autorizar el piloto.

El análisis cualitativo se aplica a las entrevistas mediante un análisis
temático, que identifica y organiza patrones de significado en los datos
\parencite{braun2006thematic}. Las notas se transcribirán y se codificarán
con categorías iniciales tomadas de la guía: procedimiento de revisión, casos
excepcionales, matrícula y afluencia, y condiciones institucionales para
tratar datos. Si aparecen temas que no caben en ellas, se agregarán como
categorías nuevas. Dentro de cada
categoría se compararán las respuestas de vigilancia y de Control Escolar
para identificar coincidencias y diferencias. Los resultados de la encuesta y
de las entrevistas sirven para el diagnóstico y los requisitos del sistema; la
respuesta a la pregunta de investigación proviene de la comparación de las
dos series de observación.

\subsection{Plan de Presentación Gráfica}
Cada resultado se presentará con la representación que mejor comunica su tipo
de dato, como resume la Tabla~\ref{tab:graficas-metodo}. Toda gráfica llevará
título, etiquetas en los ejes, unidades cuando correspondan y un párrafo de
interpretación, con el formato de figuras y tablas de APA.

\begin{table}[tbp]
	\caption{Representación Gráfica Prevista para Cada Resultado}
	\label{tab:graficas-metodo}
	\small\linespread{1}\selectfont
	\begin{tabular}{@{}P{3cm}P{6.6cm}P{5.9cm}@{}}
		\toprule
		Representación & Resultado                                                                                     & Propósito                                                          \\
		\midrule
		Barras         & Tiempo medio de validación y fila máxima promedio en la línea base y en el piloto; frecuencias de las respuestas de la encuesta; incidencias por serie. & Comparar cantidades entre condiciones o categorías.               \\[0.5em]
		Líneas         & Longitud máxima de la fila y utilización en cada intervalo, de $-30$ a $+10$ minutos, con una línea para cada serie. & Mostrar cómo cambia la fila a lo largo de la franja de entrada.    \\[0.5em]
		Circular       & Disponibilidad de celular propio, conexión y batería al llegar al plantel.                    & Mostrar proporciones de un total con pocas categorías.             \\[0.5em]
		Dispersión     & Tasa de llegadas frente a longitud máxima de la fila en cada intervalo, distinguiendo las dos series. & Explorar si la fila crece con la llegada y si esa relación cambia con el prototipo. \\[0.5em]
		Tablas         & Estadísticos descriptivos de cada serie, resultado de la prueba de hipótesis y resultados de las pruebas de banco. & Presentar valores exactos que una gráfica no muestra.              \\
		\bottomrule
	\end{tabular}
	\tablenote{Elaboración propia. Las gráficas se generarán con los mismos
		scripts del análisis.}
\end{table}
